\begin{table}[t]
\centering
\small
\resizebox{\ifdim\width>\linewidth\linewidth\else\width\fi}{!}{%%
\begin{tabular}{lllll}
\toprule
axis & levels & in the product & measure & exclusions \\
\midrule
learner & 2-4 & yes & equal / ref. / conc. & none \\
encoding & 2-3 & no: in the learner & set by the learner & none \\
split & 6 & yes & equal / ref. / conc. & none \\
decision time (availability) & 1-2 & no: second surface & never averaged & two times on the primary log, one elsewhere \\
register quality x severity & 6-10 & yes & equal / ref. / conc. & none \\
baseline (ladder rung) & 4-5 & no: the row below & equal / ref. / conc. & intercept-only: in the surface, out of every admissible set \\
baseline, admissible only & 3-4 & yes & equal / ref. / conc. & none \\
scalar instrument & 5 & yes & never averaged across & none \\
operating point (net benefit) & 31 & no: decision curve & declared threshold dist. & outside the declared range \\
\bottomrule
\end{tabular}%
}
\caption{The declared design space. One row per axis: the number of levels (a range where it differs by log), whether the axis is a FACTOR OF THE ADMISSIBLE CELL COUNT, the measure the paper places on those levels, and every level excluded by declaration with its reason. Generated from the code that walks the space. NOT EVERY ROW MULTIPLIES: only the rows marked `yes' do, and their product is learner $\times$ split $\times$ register quality $\times$ admissible baseline $\times$ scalar instrument, which is the `learners', `splits', `quality', `rungs' and `metrics' columns of Table~\ref{tab:denominator} and equals its `declared scalar' column on every pair --- the multiplication the verification harness re-does, pair by pair, from the same file this table is generated from. The rows marked `no' are not axes the study leaves unvaried; each is counted somewhere else, and the column says where. The encoding is determined by the learner --- each learner names its own encoder --- so the learner row already counts the pipeline and multiplying by the encoding would count it twice; Table~\ref{tab:axes2} separates the two where they can be separated. The decision time indexes a SECOND surface on the case-study log rather than a further factor of the first, and Section~\ref{sec:twodecision} reports it. The ladder-rung row is the admissible row plus the intercept-only baseline, so the two rows are one axis at two admissibility conventions and only the admissible one enters an admissible count. And the operating point multiplies the decision-curve surface, not the scalar one, which is why the two surfaces are counted separately in Section~\ref{sec:denominator}.}
\label{tab:axes}
\end{table}
